\documentclass[12pt]{article}
\usepackage[utf8]{inputenc}
\usepackage[T1]{fontenc}
\usepackage{hyperref}
\title{Ticuna Language Resources for K-12 Curriculum Development}
\author{Compiled by Claude Sonnet 5 (claude-sonnet-5) for Andr\'as Kornai}
\date{}
\begin{document}
\maketitle

\section{Introduction}
Ticuna (Tikuna) is spoken in Brazil, Colombia and Peru. Ticuna teachers began producing bilingual materials and a writing system in the early 1980s. The search found mainly academic studies and archive documents, not the materials themselves.

\section{Primary Educational Resources}
\begin{itemize}
\item An archive document describes Ticuna teachers creating a writing code and bilingual teaching materials. It names the myth book \emph{Tor\"u Ou\"u'\"ug\"u} and the writing manual \emph{Ngi'\~a Tana\"utchic\"unaag\"u}: \url{https://acervo.socioambiental.org/sites/default/files/documents/TCD00174.pdf}
\item A related archive document that argues for literacy in the mother tongue, taught by trained Ticuna teachers: \url{https://acervo.socioambiental.org/sites/default/files/documents/F3D00024.pdf}
\item Study of the didactic material and pedagogical project of the Ebenezer indigenous school: \url{https://bdtd.ibict.br/vufind/Record/UPM_4e1c86e76a58d8e36df7bda3ad940d88}
\end{itemize}

\section{Linguistic and Reference Materials}
\begin{itemize}
\item Soares, on translation, archives and didactic material for the Ticuna (LIAMES, 2017): \url{https://periodicos.sbu.unicamp.br/ojs/index.php/liames/article/view/8648466}
\item UFRJ dissertation on teaching Ticuna in schools in Benjamin Constant, AM: \url{https://pantheon.ufrj.br/handle/11422/25425}
\item Unicamp thesis on mathematics education for Ticuna teachers in the Alto Solim\~oes: \url{https://www.repositorio.unicamp.br/acervo/detalhe/229749}
\end{itemize}

\section{Religious and Community Resources}
None were found in search.

\section{Digital and Audio Resources}
None were found in search.

\section{Curriculum Development Context}
Ticuna teachers have produced bilingual materials since the 1980s, and schools teach the language. The materials named above should be located and checked directly. The search gave no current information on official policy or on how materials are distributed.
\end{document}